\documentclass[12pt]{article}

\usepackage{fullpage}
\usepackage{graphicx}
\usepackage{amssymb}
\usepackage{amsmath}
\usepackage[none]{hyphenat}
\usepackage{parskip}
\usepackage[spanish]{babel}
\usepackage[utf8]{inputenc}
\usepackage{hyperref}
\usepackage{fancyhdr}
\usepackage{tasks}
\usepackage{mdframed}
\usepackage{xcolor}
\usepackage{pgfplots}
\usepackage[makeroom]{cancel}
\usepackage{multicol}
\usepackage[shortlabels]{enumitem}
\usepackage{stackrel}
\usepackage{tkz-tab}
\usepackage{xpatch}
\xpatchcmd{\tkzTabLine}{$0$}{$\bullet$}{}{}

\setlength{\headheight}{10pt}
\setlength{\headsep}{10pt}
\pagestyle{fancy}
\rhead{\ayudantia \ - \alumno}
\tikzset{t style/.style={style=solid}}

\newcommand*{\mybox}[2]{\colorbox{#1!30}{\parbox{.98\linewidth}{#2}}}

\newenvironment{solucion}
{\begin{mdframed}[backgroundcolor=black!10]
		{\bf Solución:}\\
	}
	{
	\end{mdframed}
}

\newenvironment{alternativas}[1]
{\begin{multicols}{#1}
		\begin{enumerate}[a)]
		}
		{
		\end{enumerate}
	\end{multicols}
}

\newenvironment{preguntas}
{\begin{enumerate}\itemsep12pt
	}
	{
	\end{enumerate}
}

\newcommand{\ayudantia}{{\sc Ayudantía 12}}
\newcommand{\tituloayu}{Repaso I3}
\newcommand{\fecha}{4 de junio de 2019}
\newcommand{\sigla}{MAT1610}
\newcommand{\nombre}{Cálculo I}
\newcommand{\profesor}{Amal Taarabt}
\newcommand{\ano}{2019}
\newcommand{\semestre}{1}
\newcommand{\mail}{mat1610@ifcastaneda.cl}
\newcommand{\alumno}{Ignacio Castañeda - \mail}

\newcommand{\ev}{\Big|}
\newcommand{\ra}{\rightarrow}
\newcommand{\lra}{\leftrightarrow}
\newcommand{\N}{\mathbb{N}}
\newcommand{\R}{\mathbb{R}}
\newcommand{\Exp}[1]{\mathcal{E}_{#1}}
\newcommand{\List}[1]{\mathcal{L}_{#1}}
\newcommand{\EN}{\Exp{\N}}
\newcommand{\LN}{\List{\N}}
\newcommand{\comment}[1]{}
\newcommand{\lb}{\\~\\}
\newcommand{\eop}{_{\square}}
\newcommand{\hsig}{\hat{\sigma}}
\newcommand{\widesim}[2][1.5]{
	\mathrel{\overset{#2}{\scalebox{#1}[1]{$\sim$}}}
}
\newcommand{\wsim}{\widesim{}}
\newcommand{\lh}{\stackrel{L'H}{=}}

\begin{document}
\thispagestyle{empty}

\begin{minipage}{2cm}
	\includegraphics[width=2cm]{../../../../img/logo.pdf}
	\vspace{0.5cm}
\end{minipage}
\begin{minipage}{\linewidth}
	\begin{tabular}{lrl}
		{\scriptsize\sc Pontificia Universidad Catolica de Chile} & \hspace*{0.7in}Curso: &
		\sigla  - \nombre\\
		{\sc Facultad de Matemáticas}&
		Profesor: & \profesor \\
		{\sc Semestre \ano-\semestre} & Ayudante: & {Ignacio Castañeda}\\
		& {Mail:} & \texttt{\mail}
	\end{tabular}
\end{minipage}

\vspace{-10mm}
\begin{center}
	{\LARGE\bf \ayudantia}\\
	\vspace{0.1cm}
	{\tituloayu}\\
	\vspace{0.1cm}
	\fecha\\
	\vspace{0.4cm}
\end{center}

\begin{preguntas}
\item Calcule $\lim\limits_{n \ra \infty} \sqrt[]{\dfrac{1}{n^3}} + \sqrt[]{\dfrac{2}{n^3}} + \dots  + \sqrt[]{\dfrac{1}{n^2}}$
\begin{solucion}

\end{solucion}
\item Calcular
\begin{tasks}(3)
\task $\displaystyle\int_0^{\frac{1}{\sqrt[]{2}}} \dfrac{x \arcsin(x^2)}{\sqrt[]{1-x^4}}dx$
\task $\displaystyle\int_0^{\ln(2)} e^x\ \sqrt[]{e^x-1}dx$
\task $\displaystyle\int_{-3}^{4} |x-1|dx$
\end{tasks}
\begin{solucion}

\begin{enumerate}[a)]
\item 
\item 
\item 
\end{enumerate}
\end{solucion}
\item Demuestre que la función
$$F(x) = \displaystyle\int_{1-x}^{1+x} \ln(t^2)dt, \quad x \in [0, \frac{1}{2}]$$
es decreciente
\begin{solucion}

\end{solucion}
\item Determine el área limitada por la parabola $y^2=4x$ y la recta $y=x$.
\begin{solucion}
\begin{center}
			\begin{tikzpicture}
			\begin{axis}[
			axis lines = left,
			xlabel = $x$,
			ylabel = $y$,
			]
			\addplot [
			domain=0:4,  
			color=red,
			]
			{2*x^(1/2)};
			\addplot [
			domain=0:4, 
			color=blue,
			]
			{x};
			
			\end{axis}
			\end{tikzpicture}
		\end{center}
		Integramos en el eje $y$, en $[0,4]$. Para esto, usaremos las ecuaciones $x=y$ y $x=\dfrac{y^2}{4}$, por lo que el área buscada será
		$$A = \displaystyle\int_0^4 y-\dfrac{y^2}{4} dy$$
		$$A = \dfrac{16}{2} - \dfrac{64}{12} = \dfrac{8}{3}$$
\end{solucion}
\item Usando integrales, calcule el área del siguiente triangulo
	\begin{center}
		\begin{tikzpicture}
		\begin{axis}[
		axis lines = left,
		xlabel = $x$,
		ylabel = $y$,
		]
		\addplot [
		domain=0:2,  
		color=red,
		]
		{x};
		\addplot [
		domain=0:1, 
		color=red,
		]
		{-x};
		\addplot [
		domain=1:2, 
		color=red,
		]
		{3*x-4};
		
		\end{axis}
		\end{tikzpicture}
	\end{center}
\begin{solucion}
Notemos que los lados del triangulo se pueden representar como rectas. Esto es,
		$$L_1:\quad y = x$$
		$$L_2:\quad y = -x$$
		$$L_3:\quad y = 3x-4$$
		Luego, al área de triangulo será
		$$A = \displaystyle\int_0^1 (x - (-x))dx + \displaystyle\int_1^2 (x - (3x-4))dx$$
		$$A = \displaystyle\int_0^1 2xdx + \displaystyle\int_1^2 (4-2x)dx$$
		$$A = 2$$
\end{solucion}
\end{preguntas}
\end{document}